\documentclass[aspectratio=169]{beamer}

% Configuración del tema y colores UCB
\usetheme{Madrid}
\usecolortheme{default}
\definecolor{ucbblue}{RGB}{0, 51, 102}

% Ajuste de contrastes
\setbeamercolor{structure}{fg=ucbblue}
\setbeamercolor*{palette primary}{fg=white,bg=ucbblue}
\setbeamercolor*{palette secondary}{fg=white,bg=ucbblue!80}
\setbeamercolor*{palette tertiary}{fg=white,bg=ucbblue!90}
\setbeamercolor{title}{fg=white, bg=ucbblue}
\setbeamercolor{frametitle}{fg=white, bg=ucbblue}
\setbeamercolor{item}{fg=ucbblue}

% Paquetes estándar
\usepackage[utf8]{inputenc}
\usepackage[spanish]{babel}
\usepackage{amsmath, amssymb, bm}
\usepackage{graphicx}
\usepackage{booktabs}
\usepackage{tikz}
\usetikzlibrary{shapes.geometric, arrows.meta, positioning, calc}

% Información del documento
\title[Transformaciones Homogéneas SE(3)]{El Grupo Especial Euclidiano $SE(3)$ e Inversión Analítica}
\subtitle{Robótica Industrial (IMT-342) — Universidad Católica Boliviana Tarija}
\author[B. Quiroga Turdera]{M.Sc. Ing. Bernardo Quiroga Turdera}
\date{Martes 18 de Agosto de 2026}

\begin{document}

\begin{frame}
    \titlepage
\end{frame}

% Diapositiva 1
\begin{frame}{Limitaciones de $SO(3)$ y la Necesidad de $SE(3)$}
    \textbf{El Problema de las Transformaciones Afines:} \\
    Una transformación rígida general en $\mathbb{R}^3$ combina una rotación $R \in SO(3)$ y una traslación $p \in \mathbb{R}^3$:
    \begin{equation*}
        p_A = {}^{A}R_B \cdot p_B + {}^{A}p_{B,\text{org}}
    \end{equation*}
    
    \vspace{0.2cm}
    \begin{alertblock}{Obstáculo Algebraico}
        Esta ecuación \textbf{no es lineal} respecto a la posición. No puede encadenarse mediante simples multiplicaciones matriciales continuas cuando tenemos múltiples eslabones ($N$ articulaciones).
    \end{alertblock}

    \vspace{0.2cm}
    \textbf{La Solución: Espacio Proyectivo $\mathbb{P}^3$} \\
    Proyectar $\mathbb{R}^3$ a $\mathbb{P}^3$ añadiendo una coordenada escalar homogénea unitaria:
    \begin{equation*}
        \tilde{p} = \begin{bmatrix} x \\ y \\ z \\ 1 \end{bmatrix} \in \mathbb{P}^3
    \end{equation*}
\end{frame}

% Diapositiva 2
\begin{frame}{El Grupo Especial Euclidiano $SE(3)$}
    \textbf{Estructura Matricial por Bloques ($4 \times 4$):}
    \begin{equation*}
        ^{A}T_B = \begin{bmatrix} 
        \begin{array}{c|c} 
        ^{A}R_{B(3\times3)} & ^{A}p_{B(3\times1)} \\ \hline 
        \mathbf{0}_{(1\times3)} & 1 
        \end{array} 
        \end{bmatrix} = \begin{bmatrix} 
        r_{11} & r_{12} & r_{13} & p_x \\ 
        r_{21} & r_{22} & r_{23} & p_y \\ 
        r_{31} & r_{32} & r_{33} & p_z \\ 
        0 & 0 & 0 & 1 
        \end{bmatrix}
    \end{equation*}
    
    \vspace{0.4cm}
    \textbf{Desglose de Bloques Funcionales:}
    \begin{itemize}
        \item \textbf{Rotación ($3 \times 3$):} Orientación de $\{B\}$ respecto a $\{A\}$ ($R \in SO(3)$).
        \item \textbf{Traslación ($3 \times 1$):} Posición del origen de $\{B\}$ respecto al origen de $\{A\}$.
        \item \textbf{Perspectiva ($1 \times 3$):} $\mathbf{0} = [0, 0, 0]$ en robótica de cuerpo rígido.
        \item \textbf{Factor de Escala ($1 \times 1$):} Factor unitario $1$.
    \end{itemize}
\end{frame}

% Diapositiva 3
\begin{frame}{¿Qué Representa una Matriz Homogénea?}
    \renewcommand{\arraystretch}{1.5}
    \resizebox{\textwidth}{!}{
    \begin{tabular}{l l c}
    \toprule
    \textbf{Interpretación} & \textbf{Significado Físico / Aplicación} & \textbf{Ecuación Operativa} \\ \midrule
    \textbf{1. Descripción de Pose} & Almacena la ubicación y orientación completa de un eslabón. & $^{A}T_B = [^{A}\hat{x}_B, ^{A}\hat{y}_B, ^{A}\hat{z}_B, ^{A}p_B]$ \\ 
    \textbf{2. Transformación Pasiva} & Cambia la base de un punto medido en $\{B\}$ hacia el marco $\{A\}$. & $\tilde{p}_A = {}^{A}T_B \cdot \tilde{p}_B$ \\ 
    \textbf{3. Operador Activo} & Mueve y reorienta un objeto dentro del mismo marco coordenado. & $\tilde{p}' = T(\Delta p, \Delta\theta) \cdot \tilde{p}$ \\ \bottomrule
    \end{tabular}
    }
\end{frame}

% Diapositiva 4
\begin{frame}{Multiplicación Secuencial de Eslabones}
    \begin{center}
    \begin{tikzpicture}[node distance=1.8cm, auto, thick, >=stealth, scale=0.9, transform shape]
        \node[circle, draw=black, fill=gray!20] (N0) {$\{0\}$};
        \node[circle, draw=ucbblue, fill=ucbblue!10, right=of N0] (N1) {$\{1\}$};
        \node[right=of N1] (dots) {\dots};
        \node[circle, draw=ucbblue, fill=ucbblue!10, right=of dots] (N6) {$\{6\}$};
        \node[circle, draw=red!80!black, fill=red!10, right=of N6, xshift=0.5cm] (TCP) {TCP};
        
        \draw[->, ucbblue] (N0) -- node[above] {$^0T_1$} (N1);
        \draw[->, ucbblue] (N1) -- node[above] {$^1T_2$} (dots);
        \draw[->, ucbblue] (dots) -- node[above] {$^5T_6$} (N6);
        \draw[->, red!80!black] (N6) -- node[above] {$^6T_{\text{TCP}}$} (TCP);
        \draw[->, bend right=25, dashed, thick, black] (N0) to node[below] {\textbf{Pose Global:} $^0T_{\text{TCP}}$} (TCP);
    \end{tikzpicture}
    \end{center}
    
    \vspace{0.2cm}
    \textbf{Ecuación de Encadenamiento:}
    \begin{equation*}
        ^{0}T_{\text{TCP}} = {}^{0}T_{1} \cdot {}^{1}T_{2} \cdot {}^{2}T_{3} \cdot {}^{3}T_{4} \cdot {}^{4}T_{5} \cdot {}^{5}T_{6} \cdot {}^{6}T_{\text{TCP}}
    \end{equation*}
    
    \vspace{0.2cm}
    \begin{block}{Reglas Fundamentales de Multiplicación}
        \begin{itemize}
            \item \textbf{Post-multiplicación (Derecha):} Transformación respecto al marco actual móvil.
            \item \textbf{Pre-multiplicación (Izquierda):} Transformación respecto al marco inercial fijo.
        \end{itemize}
    \end{block}
\end{frame}

% Diapositiva 5
\begin{frame}{Inversión Rápida sin Eliminación Gaussiana}
    \textbf{Desarrollo Analítico:}\\
    Partiendo de la relación directa $p_A = R \cdot p_B + p \implies p_A - p = R \cdot p_B$. \\
    Multiplicando por la izquierda por $R^T$ (dado que $R^{-1} = R^T$):
    \begin{equation*}
        R^T(p_A - p) = R^T R \cdot p_B \implies p_B = R^T p_A - R^T p
    \end{equation*}
    
    En formato de bloque homogéneo $4 \times 4$:
    \begin{equation*}
        \begin{bmatrix} p_B \\ 1 \end{bmatrix} = \begin{bmatrix} R^T & -R^T p \\ \mathbf{0} & 1 \end{bmatrix} \begin{bmatrix} p_A \\ 1 \end{bmatrix}
    \end{equation*}
    
    \vspace{0.2cm}
    \begin{exampleblock}{Ecuación Fundamental de la Inversa en $SE(3)$}
        \centering $\mathbf{(^{A}T_B)^{-1} = \begin{bmatrix} R^T & -R^T p \\ \mathbf{0}_{1\times3} & 1 \end{bmatrix}}$
    \end{exampleblock}
    
    \vspace{0.1cm}
    \small \textit{\textbf{Costo Computacional:} La inversión analítica requiere $\approx 10$ operaciones. La eliminación numérica genérica (Gauss/LU) supera las $60$ operaciones acumulando error de truncamiento.}
\end{frame}

% Diapositiva 6
\begin{frame}{Caso Práctico PFI — Calibración de la Garra Neumática}
    \textbf{Problema Real (ABB IRB 120):} El controlador conoce $^{0}T_6$ (brida). La pinza añade un desplazamiento de $L = 120\text{ mm}$ en $Z_6$ y sin desfase de rotación.
    
    \vspace{0.2cm}
    \begin{columns}
        \begin{column}{0.6\textwidth}
            \textbf{Matriz de Herramienta ($^{6}T_{\text{TCP}}$):}
            \begin{equation*}
                ^{6}T_{\text{TCP}} = \begin{bmatrix} 
                1 & 0 & 0 & 0 \\ 
                0 & 1 & 0 & 0 \\ 
                0 & 0 & 1 & 0.120 \\ 
                0 & 0 & 0 & 1 
                \end{bmatrix}
            \end{equation*}
            
            \textbf{Cálculo de la Pose Real:}
            \begin{equation*}
                ^{0}T_{\text{TCP}} = {}^{0}T_6(q) \cdot {}^{6}T_{\text{TCP}}
            \end{equation*}
        \end{column}
        \begin{column}{0.4\textwidth}
            \centering
            \begin{tikzpicture}[scale=1, axis/.style={->, thick, >=stealth}]
                % Brida {6}
                \draw[thick, fill=gray!30] (-0.5,-0.2) rectangle (0.5,0.2);
                \draw[axis, blue] (0,0) -- (1,0) node[right]{$X_6$};
                \draw[axis, black] (0,0) -- (0,1) node[left]{$Z_6$};
                
                % Pinza TCP
                \draw[thick, fill=red!10] (-0.1, 0.2) -- (-0.1, 1.8) -- (0.1, 1.8) -- (0.1, 0.2) -- cycle;
                \draw[thick, fill=red!30] (-0.3, 1.8) rectangle (-0.1, 2.2);
                \draw[thick, fill=red!30] (0.1, 1.8) rectangle (0.3, 2.2);
                
                % TCP Frame
                \draw[axis, blue] (0,2) -- (1,2) node[right]{$X_{\text{TCP}}$};
                \draw[axis, black] (0,2) -- (0,3) node[left]{$Z_{\text{TCP}}$};
                
                % Cota
                \draw[<->, dashed] (-0.6, 0) -- node[left]{$120\text{mm}$} (-0.6, 2);
            \end{tikzpicture}
        \end{column}
    \end{columns}
    
    \vspace{0.3cm}
    \small \textbf{Extracción Posicional:} Las 3 primeras filas de la cuarta columna de $^{0}T_{\text{TCP}}$ indican la posición métrica $(x,y,z)$ de cierre de las mordazas.
\end{frame}

\end{document}